%% 5. Veille technologique  (FACULTATIF mais valorisé)
%% COMPTE dans la règle "Veille + Travaux >= 50 % du rapport".
\chapter{Veille technologique}\label{ch:veille}

% Ce chapitre décrit la veille technique CONTINUE menée tout au long de
% l'alternance. La veille spécifique à chaque projet (choix d'architecture,
% évaluation d'OCR, pistes d'augmentation de données) est désormais présentée
% directement dans le chapitre Travaux (\ref{ch:travaux}).

Tout au long de mon alternance, le suivi des évolutions technologiques a constitué
une activité \textbf{continue}, menée en parallèle des missions proprement dites. Le
domaine des grands modèles de langage et de leur mise en production évoluant à un
rythme particulièrement soutenu, cette veille était indispensable pour maintenir à
l'état de l'art nos outils, au premier rang desquels le \textbf{Docprocessor} (voir
chap.~\ref{ch:env}), et pour anticiper les opportunités d'amélioration.

\section{Suivi des nouvelles versions de vLLM}

\textbf{vLLM} est le moteur d'inférence et de service open source sur lequel
reposent nos déploiements de grands modèles de langage. Son développement étant très
actif, j'ai assuré un \textbf{suivi régulier de ses nouvelles versions} afin
d'identifier les apports susceptibles de bénéficier à nos déploiements.
% À COMPLÉTER : versions suivies, apports concrets (gestion mémoire / PagedAttention,
% débit, sortie structurée, quantification...), et impact sur nos déploiements.
d'une mise à jour à une autre, les performances de mes modèles pouvaient variés (genre il bouclait)

\section{Serveurs de modèles et solutions de service}

% À COMPLÉTER : panorama des serveurs de modèles et solutions de service suivis et
% comparés (vLLM et alternatives), critères de comparaison (débit, latence, coût,
% fonctionnalités supportées) et ce qui a orienté nos choix d'infrastructure.
j'ai du commencé à veilelr sur docprocessor par grafana car à un moment j'ai travaillé sur une optimisation en temps de la pipeline en passant de deux requêtes à une requêtes qui faisait tous.
sauf que du coup fallait vraiment monitorer que tout fonctionne.



\section{Fonctionnalités et tendances émergentes}
pour entraîner les vlm, j'ai utilisé unsloth qui une surcouche. ce moteur est basé sur les libraries transformers, trl etc donc c'est assez compliqué de faire quelque chose qui respecte toutes les mise à jour
sachant que le monde de l'ia générative va très très vite. du coup c'est vraiment du code auto généré et à base de monkey-patch tout le temps donc pareil d'une version à une autre beaucoup de choses pouvaient changer.


% À COMPLÉTER : fonctionnalités « tendance » suivies (sortie structurée / contrainte
% par schéma, nouveaux modèles, OCR de nouvelle génération...) et veille sur les
% usages applicables à nos produits de vérification d'identité.

